The prior scoping review of cybersickness EEG research by Chang et al.~\cite{changBrainActivityCybersickness2023} reported delta-band increases, beta-band decreases, and inconsistent alpha-band findings; the present corpus refines this picture in two respects.
The delta- and theta-band increases reproduce in the current corpus (16 of 30 studies for delta; 20 of 35 for theta).
The reported alpha-band inconsistency partly resolves once VR display hardware is accounted for: head-mounted-display and screen-based studies differ with large effects on alpha (Cliff's $\delta = 0.875$, $p = 0.004$) and beta ($\delta = 0.664$, $p = 0.015$).

\subsection{Objective Markers and the Limits of Self-Report}
VIMS has no diagnostic biomarker, its ground truth is self-report.
It is defined by its symptom profile, nausea, disorientation, and oculomotor disturbance, which are private experiences that cannot be externally verified \cite{kennedySimulatorSicknessQuestionnaire1993, stanneyCybersicknessNotSimulator1997}; any candidate objective marker must therefore be calibrated against self-report rather than replacing it.
Functional cognitive consequences of this state are documented in 2 corpus studies: working-memory decrements scaling with reported severity \cite{mimnaughVirtualRealitySickness2023} and inhibitory-control demand reflected in N2/P3 components \cite{wuInhibitionrelatedN2P32020}, confirming that the state extends beyond subjective discomfort and carries measurable functional costs.
The EEG and fNIRS correlates established in \autoref{sec:correlates} are objective physiological measurements, but each was validated against self-reported severity and therefore presupposes the measure it would replace.
Among the candidate objective markers, postural sway has a direct theoretical anchor in the postural-instability account, which predicts that pre-symptomatic sway covaries with subsequent severity \cite{riccioEcologicalTheoryMotion1991, stoffregenPosturalInstabilityPrecedes1998}; VR-specific reports document moderate coupling \cite{palmisanoSpontaneousPosturalSway2014, dennisonImprovingMotionSickness2019, teixeiraInvestigatingRelativeContributions2024}.
Pupillometry is a second candidate \cite{kourtesisCybersicknessVirtualReality2023, kourtesisCybersicknessCognitionMotor2023}, subject to arousal and cognitive-load confounds.
Of the 92 corpus studies, 16 already co-record force-plate, balance-board, eye-tracker, or pupillometry alongside EEG; treating these signals as candidate outcome variables rather than covariates is a prerequisite for validating them as objective markers.

\subsection{Methodological Challenges and Critical Perspective}
% To highlight: VR headsets, resolutions, high accuracy, young subjects, small samples

While the reviewed literature demonstrates progress in applying neuroimaging and neurostimulation methods to cybersickness research, a critical appraisal reveals substantial methodological challenges.
These issues temper the interpretation of many findings and present considerable barriers to the generalization, replication, and practical application of the reported results.
This section will highlight four main challenges identified in the literature (1) optimistic accuracy estimates, (2) the lack of standardized cybersickness thresholds, (3) protocol disparities and (4) conceptual variability.

A prominent trend in the classification literature is an "accuracy arms race," with a proliferation of studies reporting exceptionally high performance.
Numerous papers claim binary classification accuracies exceeding 95\%, and in some cases, approaching 99\% \cite{huaDetectionVirtualReality2024, shenClassificationVisuallyInduced2024, huaEEGClassificationModel2024, dasdemirLocomotionTechniquesEEG2023, huaDualpathwayEEGModel2025}.
However, these impressive figures are often a product of experimental designs that limit their generalizability.
Many are achieved using within-subject validation on small, homogeneous participant pools, which may not translate to new, unseen users.
Furthermore, the choice of data used for classification also plays a role. For example, some high-performing models compare resting-state EEG data collected before and after VR exposure \cite{huaEEGClassificationModel2024, chaiEEGStudyVirtual2023}, a paradigm that, while useful for identifying state changes, is ill-suited for the real-time detection required by adaptive systems.

The foundation of any supervised learning model is the quality of its "ground truth" labels, and in cybersickness research, this foundation is ill-defined.
The vast majority of predictive models are trained on subjective data, most commonly post-exposure SSQ scores or in-task ratings from scales like the FMS \cite{rosannePreprocessNotPreprocess2025}.
The SSQ is retrospective and coarse, while continuous ratings can lead the model to classify signals related to the motor action of reporting.
A more fundamental issue is the lack of standardized thresholds for defining sickness states.
Individual studies frequently devise their own arbitrary cut-offs to convert continuous SSQ scores into binary ("sick" vs. "not sick") or multi-level ("mild," "moderate," "severe") labels.
This means that the very definition of the target variable is inconsistent across the literature, severely undermining the comparability of different models and their reported accuracies.

This lack of standardization extends to many aspects of experimental design, creating a profound heterogeneity that hinders scientific synthesis.
First, there is significant variability in display technology and the nature of the stimulus.
Experiences in fully immersive HMDs are not equivalent to those on large-screen monitors or in CAVE-like projection systems \cite{limTestretestReliabilityVirtual2021, zhangAnalysisMotionSickness2020}, or smartphone based VR systems \cite{molefiVibromotorReprocessingTherapy2022}, yet findings are often discussed interchangeably.
Similarly, passive viewing of a 360-degree video \cite{suwandiExploringImpactMotion2024} is fundamentally different from active navigation in an interactive game \cite{cortesEEGbasedExperimentVR2023a}.
Our meta-analysis supports these concerns, revealing statistically significant differences in neural response patterns between different hardware platforms and virtual environment types (see \autoref{subsubsec:confounding}).
Second, experimental protocols are inconsistent, with wide variance in exposure duration, the inclusion and length of rest periods, and the definition of a valid baseline.
A particularly problematic practice is the use of an "eyes-closed" baseline, which artificially inflates alpha power and directly confounds the interpretation of alpha suppression, one of the most commonly reported neural markers of cybersickness.
Third, the recording itself varies, ranging from research-grade, high-density EEG systems to low-cost, consumer-grade devices with fewer channels and potentially lower signal quality, which may impact the reliability of the derived results \cite{liuPilotStudyEEGBased2020}.
Additionally, reporting quality across studies demonstrates substantial deficiencies in critical methodological parameters.
For example, EEG reference electrode configuration is reported in only 49.4\% of studies, and notch filter frequency in 22.9\%, hampering reproducibility and meta-analytic synthesis.
A quantitative analysis of the included literature reveals a reliance on low-density EEG systems, with 82.9\% of studies employing 32 or fewer electrodes, and 21.4\% employing 10 or fewer electrodes.
While such systems are appropriate for applied BCI paradigms that target well-characterized and localized neural signals, their use in the context of cybersickness research presents a significant methodological limitation.
The neural underpinnings of cybersickness are not yet fully elucidated, and the research objective for many studies remains fundamentally exploratory.
Low-density electrode montages have insufficient spatial resolution for reliable source localization and are blind to activity originating from large, un-sampled cortical areas.
This practice, therefore, risks generating an incomplete understanding of cybersickness neurophysiology, potentially overlooking critical contributions from distributed neural networks and biasing the literature towards the few cortical sites that are consistently measured.
Compounding these protocol issues, the literature is dominated by small and demographically narrow samples (typically young, predominantly male university students), which further limits the statistical power and external validity of the reported findings \cite{jangEstimatingObjectiveEEG2022, molefiVibromotorReprocessingTherapy2022}.

Finally, the field exhibits conceptual variability in both the use of theory and the use of terminology.
Many studies cite an etiological theory of cybersickness in their introduction or methods, but comparatively few return to that framework when interpreting their neurophysiological findings.
The consequence is interpretive ambiguity: a neural feature observed in a given paradigm could reflect sensory conflict, postural control, expectation-related vection, display optics, or some combination thereof, depending on the exposure context, and without an explicit theoretical commitment the literature cannot adjudicate between these possibilities.
Terminological practice compounds the problem.
Motion sickness, VIMS, simulator sickness, and cybersickness are related but non-identical constructs, each tied to a different exposure context (see \autoref{sec:theoretical}), yet they are frequently used interchangeably across studies \cite{kennedyResearchVisuallyInduced2010, chaMotionSicknessDiagnostic, stanneyCybersicknessNotSimulator1997, rebenitschReviewCybersicknessApplications2016}.
A similar conflation arises between sickness symptoms and the perceptual phenomenon of vection, despite the latter being absent from the SSQ and from the Bárány diagnostic criteria \cite{hettingerVectionSimulatorSickness1990, kennedySimulatorSicknessQuestionnaire1993, chaMotionSicknessDiagnostic}.
A further consequence of this conceptual looseness is the tendency to label any statistically significant feature a "neuromarker" without establishing its specificity, sensitivity, or reliability across contexts.

Taken together, this landscape of methodological and conceptual variability poses a significant threat to replicability and impedes the cumulative progress of the field.

\subsection{Future Directions}\label{sec:future}
The methodological challenges documented in \autoref{sec:discussion} translate into six concrete directions for future work: standardised reporting and open data, causal mechanisms, generalisation and ecological validation, sex- and gender-stratified design, objective behavioural markers, and longitudinal characterisation of adaptation and recovery.

First, standardised reporting represents a gap documented across the corpus.
For classification and prediction studies, the minimum items follow from the gaps documented in \autoref{sec:discussion}: a discriminative metric (AUC or balanced accuracy) with 95\% confidence interval, sensitivity and specificity at a stated operating point, the numeric label threshold, a subject-respecting validation scheme, and comparison against a trivial and a non-neural feature baseline.
For studies that quantify in-exposure neural change, a baseline matched to the exposure condition on visual input, oculomotor state, and vigilance is required for reproducibility, with the choice and its rationale reported; the available options are discussed in \autoref{sec:discussion}.
For EEG, the reference electrode and notch-filter frequency are required for reproducibility and meta-analytic re-pooling.
For subjective sickness measurement, reporting the questionnaire instrument, its administration time-point relative to exposure end, and the numeric cutoff used to assign class labels enables cross-study comparability; where continuous in-task ratings are used, the response modality and any associated motor artefact exclusion strategy are necessary for reproducibility.
For stimulus and display, the device model, field-of-view, and frame rate alongside the stimulus class (passive or active, recorded or interactive); adoption of a small number of canonical reference paradigms, at minimum a passive visual-flow exposure \cite{suwandiExploringImpactMotion2024} and an active first-person navigation task \cite{cortesEEGbasedExperimentVR2023a}, would allow direct cross-laboratory comparison without the confound of idiosyncratic stimulus design.

These reporting standards are a precondition for, but not a substitute for, shared infrastructure: community-curated benchmark datasets with multimodal recordings (EEG, EOG, ECG, head motion, sway, eye tracking), well-defined subjective labels, and coverage of both passive and active paradigms.
A small number of corpus papers have already released public datasets \cite{demirelSubjectivityContinuousCybersickness2025, kimDeepCybersicknessPredictor2019, dasdemirClassificationEmotionalImmersive2023}; releasing analysis code and pre-trained models alongside the data would further support replication and cumulative analysis.

Second, the causal status of documented neural associations remains largely unexamined.
The literature has documented associations between specific neural patterns and VIMS across multiple studies, but whether these associations are causal has not been established.
Non-invasive brain stimulation offers a direct test: transcranial alternating-current stimulation can target specific oscillations (e.g., alpha or theta), and existing corpus studies that combine stimulation with VR exposure \cite{takeuchiModulationExcitabilityTemporoparietal2018, bergerInfluencePlaceboTDCS2024} provide a methodological template for asking whether modulating a candidate neuromarker changes sickness severity.
Coupled with computational models of multisensory integration and predictive coding \cite{nurnbergerMismatchVisualVestibularInformation2021}, interventional designs can begin to distinguish neural features that drive the VIMS experience from features that merely accompany it.

Third, generalisation and ecological validation remain open challenges for models built for VIMS.
In the 4 corpus studies that report both within- and cross-subject accuracy, cross-subject performance is lower in each case, and only 1 of 38 classification studies reports a leave-one-subject-out evaluation; transfer-learning and domain-adaptation approaches that explicitly model inter-subject and inter-session variability address a current gap if brain-directed methods are to operate without per-user calibration.
Validation beyond the laboratory is a further open direction: at-home, naturalistic paradigms such as the in-home gaming setup of \cite{moinnereauInstrumentingVirtualReality2022} address the question of whether laboratory neuromarkers survive the noise of consumer-grade hardware and uncontrolled environments.

Fourth, sex and gender represent an underexplored factor in the corpus.
Of the 92 corpus studies, 11 used male-only samples, 15 did not report sex composition, and only 3 (all from a single laboratory) reported higher cybersickness in women \cite{bergerSexDifferencesUser2021, bergerInfluencePlaceboTDCS2024, bergerManipulatingCybersicknessVirtual2025}.
The broader literature on sex differences in VIMS susceptibility is mixed: \cite{kellyGenderDifferencesCybersickness2023} argues that reported effects are modest in magnitude and plausibly mediated by interpupillary-distance fit, prior gaming experience, and questionnaire framing; \cite{kourtesisCybersicknessCognitionMotor2023} reports that male--female differences disappear when prior gaming experience is included as a covariate; and \cite{lawsonAttentionMilitaryCommercial2023} caution that the claim of greater female susceptibility is overstated in the simulation literature.
Taken together, the 3 corpus reports and the broader literature are consistent with a modest, variable effect whose magnitude and moderators cannot be resolved at individual-study scale, and whose apparent direction may partly reflect gaming-experience differences and questionnaire framing rather than a stable sex difference in susceptibility.
Reporting sex-stratified outcomes (band-power changes, connectivity, classification metrics), even when single-study samples are underpowered for inferential testing, would enable subsequent meta-analytic re-pooling.
Logging interpupillary-distance setting and gaming experience as covariates would address both the sex confound and the hardware effects documented in \autoref{sec:confounders}.
Adequately powered designs, such as pooled multi-site cohorts or longitudinal within-subject designs, are the precondition for resolving the magnitude and moderators of a sex effect on VIMS.

Fifth, objective behavioural markers represent an underdeveloped line of investigation.
The candidate markers introduced in \autoref{sec:discussion} are already collected by 16 of the 92 corpus studies, but typically as covariates rather than as candidate ground truth.
Postural sway is the strongest theoretically anchored candidate: reporting a coupling estimate between pre-symptomatic sway dynamics and subsequent self-reported sickness, alongside the protocol metadata that this estimate depends on (stance, surface, sampling rate), would anchor the construct across laboratories.
Pupillometry can supplement self-report subject to arousal and cognitive-load confounds \cite{kourtesisCybersicknessVirtualReality2023}, and requires controls on ambient luminance and on task demand.
Heart-rate variability, electrodermal activity, and head-motion dynamics are further candidates whose corpus-level predictive value is not yet quantified \cite{islamAutomaticDetectionPrediction2020, tianWhoSaysYou2024, yangPredictionDetectionVirtual2023}; multi-modal integration that fuses behavioural, autonomic, and neural signals offers a more holistic state estimate than any single modality.
Treating these markers as outcome variables, rather than as covariates, is the precondition for shifting VIMS research toward objective ground truth.

Finally, longitudinal designs represent a gap in the current corpus.
The corpus is dominated by single-session experiments, leaving the neurodynamics by which users adapt to VR (``getting one's VR legs'') largely uncharacterised, and the single corpus recovery study \cite{wooRecoveryTimeVR2023} is insufficient to support strong claims about residual neural state.
Multi-session within-subject designs would enable tracking of how the EEG signatures identified in this review evolve with repeated exposure, and mapping the dissipation of post-exposure signatures onto behaviour.

